\chapter{Интеллектуальные сетевого оборудования функции управляемого сетевого оборудования}
\label{ch:chap5}

Перед началом выполнения первой части лабораторной работы была выполнена инициализация виртуального окружения лабороторной работы с помощью виртуального эмулятора PNETLab.

Цель данного этапа - изучить принципы работы управляемого сетевого оборудования, смоделировать атаки канального уровня и освоить базовые механизмы защиты локальной сети, такие как Port Security, DHCP Snooping и Dynamic ARP Inspection.

\section{Подготовка к выполнению}
\begin{comment}
\begin{figure}[ht!]
    \centering
    \includegraphics[width=0.8\textwidth]{screenshots/1_1.jpg}
    \caption{Установка рабочего окружения}
    \label{fig:1_1}
\end{figure}
\clearpage
\begin{figure}[ht!]
    \centering
    \includegraphics[width=0.8\textwidth]{screenshots/1_2.jpg}
    \caption{Вход в рабочее окружение}
    \label{fig:1_2}
\end{figure}
\clearpage
\begin{figure}[ht!]
    \centering
    \includegraphics[width=0.8\textwidth]{screenshots/1_3.jpg}
    \caption{Создание новой лабороторной работы}
    \label{fig:1_3}
\end{figure}
\clearpage
\end{comment}

\section{Конфигурация узлов локальной сети}
\begin{figure}[ht!]
    \centering
    \includegraphics[width=1.0\textwidth]{screenshots/1_1.jpg}
    \caption{Воспроизведённая топология локальной сети}
    \label{fig:1_0}
\end{figure}
\clearpage
\begin{figure}[ht!]
    \centering
    \includegraphics[width=1.0\textwidth]{screenshots/1_2.jpg}
    \caption{Установка ip-адреса локального узла Kali}
    \label{fig:1_2}
\end{figure}
\clearpage
\begin{figure}[ht!]
    \centering
    \includegraphics[width=1.0\textwidth]{screenshots/1_4.jpg}
    \caption{Установка ip-адреса <<обычного>> узла}
    \label{fig:1_4}
\end{figure}
\clearpage
\begin{figure}[ht!]
    \centering
    \includegraphics[width=1.0\textwidth]{screenshots/1_5.jpg}
    \caption{Проверка доступности узла <<жертвы>> из узла <<атакующего>>}
    \label{fig:1_5}
\end{figure}
\clearpage

\section{Воспроизведение атаки MAC-spoofing}
\begin{figure}[ht!]
    \centering
    \includegraphics[width=1.0\textwidth]{screenshots/1_6.jpg}
    \caption{Получение MAC-адреса узла <<жертвы>>}
    \label{fig:1_6}
\end{figure}
\clearpage
\begin{figure}[ht!]
    \centering
    \includegraphics[width=1.0\textwidth]{screenshots/1_7.jpg}
    \caption{Подмена MAC-адреса узла <<атакующего>> на MAC-адрес узла <<жертвы>>}
    \label{fig:1_7}
\end{figure}
\clearpage
\begin{figure}[ht!]
    \centering
    \includegraphics[width=1.0\textwidth]{screenshots/1_8.jpg}
    \caption{Подмена MAC-адреса узла <<атакующего>> на MAC-адрес узла <<жертвы>>}
    \label{fig:1_8}
\end{figure}
\clearpage
\begin{figure}[ht!]
    \centering
    \includegraphics[width=1.0\textwidth]{screenshots/1_9.jpg}
    \caption{Создание траффика между <<обычным>> узлом и узлом <<жертвой>>}
    \label{fig:1_9}
\end{figure}
\clearpage
\begin{figure}[ht!]
    \centering
    \includegraphics[width=1.0\textwidth]{screenshots/1_10.jpg}
    \caption{Анализ ранее созданного траффика с помощью утилиты Wireshark}
    \label{fig:1_10}
\end{figure}
\clearpage
\begin{figure}[ht!]
    \centering
    \includegraphics[width=1.0\textwidth]{screenshots/1_11.jpg}
    \caption{Создание обманывающего траффика между узлом <<атакующим>> и узлом <<жертвой>>}
    \label{fig:1_11}
\end{figure}
\clearpage
\begin{figure}[ht!]
    \centering
    \includegraphics[width=1.0\textwidth]{screenshots/1_12.jpg}
    \caption{Перезаписи таблицы коммутатора Switch2}
    \label{fig:1_12}
\end{figure}
\clearpage
\begin{figure}[ht!]
    \centering
    \includegraphics[width=1.0\textwidth]{screenshots/1_13.jpg}
    \caption{Отбросы кадров на <<обычном>> узле из-за перезаписи MAC-адреса <<атакующим>> узлом}
    \label{fig:1_13}
\end{figure}
\clearpage

\section{Выводы по главе}

\endinput
